% ============================================================
%  CHAPITRE 2 — État de l'art
%  À inclure dans le document principal avec :
%    \input{chapitre2}   ou   \include{chapitre2}
%
%  Packages supposés chargés dans le main :
%    \usepackage[french]{babel}
%    \usepackage[utf8]{inputenc}
%    \usepackage[T1]{fontenc}
%    \usepackage{booktabs}
%    \usepackage{array}
%    \usepackage{tabularx}
%    \usepackage{xcolor}
%    \usepackage{hyperref}
%    \usepackage{geometry}
%    \usepackage{enumitem}
%    \usepackage{parskip}       % ou gestion manuelle des espacements
% ============================================================

\chapter{État de l'art}
\label{chap:etat-de-lart}

Ce chapitre situe le projet BCaaS (\textit{Business Core as a Service}) dans son paysage
scientifique et technique. Il ne s'agit pas simplement de lister des concepts existants,
mais de montrer comment chacun d'eux contribue à la conception de BCaaS et en quoi leur
combinaison constitue une approche originale. Nous passerons successivement en revue les
modèles en couches issus des réseaux informatiques, les stratégies d'architecture
multi-tenant, le paradigme \textit{as-a-Service} et son extension au niveau métier, les
principes de l'architecture hexagonale et du \textit{Domain-Driven Design}, avant de
conclure en positionnant la contribution spécifique de BCaaS par rapport à cet état de
l'art.

% ------------------------------------------------------------
\section{Modèles en couches : OSI et TCP/IP}
\label{sec:osi-tcpip}

% ---- 2.1.1 -------------------------------------------------
\subsection{Le modèle OSI}
\label{subsec:osi}

Le modèle OSI (\textit{Open Systems Interconnection}), normalisé par l'ISO sous la
référence ISO/IEC~7498-1~\cite{iso7498}, est un cadre de référence conceptuel qui décompose
la communication entre deux systèmes informatiques en sept couches aux responsabilités
strictement délimitées. Chaque couche remplit une fonction précise et ne communique
directement qu'avec les couches immédiatement adjacentes, sans connaître les détails de
fonctionnement des autres.

Les sept couches, de la plus basse à la plus haute, sont les suivantes :

\begin{description}[leftmargin=2.5cm, labelwidth=2.2cm, style=nextline]
  \item[\textbf{Couche~1 Physique}]
    Transmission des bits bruts sur le médium (câble, fibre, ondes radio). Elle définit
    les niveaux de tension, les fréquences et les connecteurs physiques.

  \item[\textbf{Couche~2 Liaison}]
    Organisation des bits en trames, détection (voire correction) d'erreurs de
    transmission, et contrôle d'accès au médium partagé (protocoles Ethernet, Wi-Fi).

  \item[\textbf{Couche~3 Réseau}]
    Adressage logique (adresses IP) et routage des paquets entre réseaux différents.
    C'est la couche qui permet à un paquet de traverser plusieurs réseaux intermédiaires
    pour atteindre sa destination.

  \item[\textbf{Couche~4 Transport}]
    Communication de bout en bout entre processus, avec ou sans garantie de fiabilité.
    TCP assure la livraison ordonnée et sans perte~; UDP privilégie la rapidité au
    détriment de cette garantie.

  \item[\textbf{Couche~5 Session}]
    Établissement, maintien et terminaison des sessions de dialogue entre deux
    applications.

  \item[\textbf{Couche~6 Présentation}]
    Mise en forme des données (encodage des caractères, compression, chiffrement) afin
    que les applications puissent se comprendre indépendamment de leurs représentations
    internes.

  \item[\textbf{Couche~7 Application}]
    Interface directe avec les logiciels utilisateurs. C'est ici que résident les
    protocoles applicatifs tels que HTTP, SMTP, FTP ou DNS.
\end{description}

Le principe central du modèle OSI est l'\textbf{encapsulation} : lorsqu'une donnée
descend la pile d'émission, chaque couche y ajoute un en-tête (et parfois un pied de
trame) contenant les informations de contrôle nécessaires à son homologue sur la machine
distante. À réception, la pile est parcourue en sens inverse~: chaque couche lit et retire
son en-tête avant de passer les données à la couche supérieure. Ce mécanisme garantit une
séparation complète des responsabilités~: chaque couche ignore le contenu des données
qu'elle transporte et ne se préoccupe que de sa propre fonction.

\begin{table}[ht]
  \centering
  \caption{Les sept couches OSI et leurs analogies dans BCaaS}
  \label{tab:osi-bcaas}
  \renewcommand{\arraystretch}{1.3}
  \begin{tabularx}{\textwidth}{>{\centering\bfseries}p{1.2cm}
                               >{\bfseries}p{2.4cm}
                               >{\raggedright\arraybackslash}X
                               >{\raggedright\arraybackslash}X}
    \toprule
    \textbf{N°} & \textbf{Couche} & \textbf{Rôle principal} & \textbf{Analogie BCaaS} \\
    \midrule
    7 & Application   & Interface avec l'utilisateur final                    & API métier exposée aux clients \\
    6 & Présentation  & Encodage, chiffrement, compression                    & Sérialisation JSON/XML, tokens JWT \\
    5 & Session       & Gestion des sessions de communication                 & Gestion des sessions utilisateur \\
    4 & Transport     & Fiabilité, contrôle de flux (TCP/UDP)                 & File d'attente, retry, transactions \\
    3 & Réseau        & Adressage et routage inter-réseaux                    & Routage multi-tenant, isolation VLAN \\
    2 & Liaison       & Détection d'erreurs, accès au médium                  & Validation des données, contraintes DB \\
    1 & Physique      & Transmission binaire sur le médium                    & Infrastructure serveur, cloud \\
    \bottomrule
  \end{tabularx}
\end{table}

% ---- 2.1.2 -------------------------------------------------
\subsection{La suite TCP/IP}
\label{subsec:tcpip}

La suite de protocoles TCP/IP, développée par le Département de la Défense américain dans
les années~1970 et standardisée au sein des RFC (\textit{Request For Comments}), adopte
une approche plus pragmatique en condensant les sept couches OSI en quatre niveaux~:
accès réseau (couches~1-2), Internet (couche~3), transport (couche~4) et application
(couches~5-7). C'est ce modèle qui régit le fonctionnement d'Internet depuis son origine.

TCP/IP conserve les mêmes idées directrices que le modèle OSI — adressage, routage,
fiabilité, contrôle de flux — mais en simplifiant leur mise en œuvre. Son adoption
universelle a été rendue possible par sa flexibilité~: l'adressage IP, notamment avec
IPv6, permet d'identifier de manière unique chaque appareil connecté à l'échelle mondiale.

% ---- 2.1.3 -------------------------------------------------
\subsection{Intérêt méthodologique pour BCaaS}
\label{subsec:osi-bcaas}

Au-delà de leur application aux réseaux informatiques, ces modèles présentent un intérêt
méthodologique fondamental pour la conception logicielle. Trois principes sont
particulièrement transposables~:

\begin{description}[leftmargin=0.5cm]
  \item[\textbf{Séparation des préoccupations par couches} (\textit{Separation of
    Concerns})]
    Chaque module ou composant logiciel est responsable d'une et d'une seule fonction. Ce
    principe réduit la complexité, facilite les tests unitaires et rend les composants
    remplaçables indépendamment.

  \item[\textbf{Encapsulation progressive du contexte}]
    Chaque couche enrichit les données d'informations supplémentaires (métadonnées,
    identifiants, contexte de sécurité) sans que les couches adjacentes n'aient à
    connaître ces détails. En logiciel, cela correspond à l'enrichissement d'un objet
    métier à mesure qu'il traverse les différentes couches d'une application (validation,
    transformation, persistance).

  \item[\textbf{Adressage explicite des destinataires}]
    Dans TCP/IP, une adresse IP et un port identifient sans ambiguïté la destination d'un
    paquet. En architecture multi-tenant, l'identifiant du tenant (\texttt{tenant\_id})
    joue un rôle analogue~: il adresse chaque requête vers le bon contexte de données et
    de règles métier.
\end{description}

C'est cette transposition que BCaaS exploite explicitement~: la pile applicative métier
est conçue comme une pile protocolaire, avec des couches aux responsabilités distinctes,
des mécanismes d'encapsulation du contexte (tenant, utilisateur, droits) et un adressage
explicite par \texttt{tenant\_id}.

% ------------------------------------------------------------
\section{Architectures multi-tenant}
\label{sec:multi-tenant}

% ---- 2.2.1 -------------------------------------------------
\subsection{Définition et enjeux}
\label{subsec:multitenant-def}

Une application est dite \textbf{multi-tenant} (multi-locataire) lorsqu'une seule instance
logicielle sert simultanément plusieurs organisations ou clients distincts (les
\textit{tenants} ) tout en garantissant l'isolation totale de leurs données et
configurations. Ce modèle s'oppose au modèle multi-instance, où chaque client dispose de
sa propre installation du logiciel~\cite{multitenant}.

Les enjeux du multi-tenant sont multiples et parfois contradictoires~:

\begin{description}[leftmargin=0.5cm]
  \item[\textbf{Isolation des données}]
    Un tenant ne doit jamais pouvoir accéder, même accidentellement, aux données d'un
    autre. Cette exigence est à la fois technique (partitionnement physique ou logique) et
    réglementaire (RGPD, normes sectorielles).

  \item[\textbf{Performance et scalabilité}]
    Le partage d'infrastructure peut introduire des effets de voisinage
    (\textit{noisy neighbor problem}) où un tenant très actif dégrade les performances
    des autres.

  \item[\textbf{Personnalisation}]
    Chaque tenant peut avoir des configurations, des règles métier ou des interfaces
    différentes, tout en partageant le même code de base.

  \item[\textbf{Opérabilité}]
    Les mises à jour, les sauvegardes et la supervision doivent être gérées de manière
    centralisée pour tous les tenants, sans interruption de service.
\end{description}

% ---- 2.2.2 -------------------------------------------------
\subsection{Trois stratégies d'isolation}
\label{subsec:isolation-strategies}

La littérature académique et les pratiques de l'industrie distinguent classiquement trois
niveaux d'isolation, chacun avec ses compromis en termes de sécurité, de performance et
de coût opérationnel~\cite{strategies}.

\begin{table}[ht]
  \centering
  \caption{Comparaison des stratégies d'isolation multi-tenant}
  \label{tab:multitenant}
  \renewcommand{\arraystretch}{1.3}
  \begin{tabularx}{\textwidth}{>{\raggedright\arraybackslash}X
                               >{\centering\arraybackslash}p{2.2cm}
                               >{\centering\arraybackslash}p{2.4cm}
                               >{\raggedright\arraybackslash}X}
    \toprule
    \textbf{Stratégie} & \textbf{Isolation} & \textbf{Performance} & \textbf{Complexité / Coût} \\
    \midrule
    Colonne \texttt{tenant\_id} (schéma partagé)
      & Logique   & Bonne (index)  & Faible idéal pour un MVP \\
    Schéma PostgreSQL séparé
      & Physique  & Très bonne     & Moyenne \\
    Base de données séparée
      & Totale    & Variable       & Élevée \\
    \bottomrule
  \end{tabularx}
\end{table}

\subsubsection*{a) Schéma partagé avec colonne \texttt{tenant\_id}}

C'est la stratégie la plus simple. Toutes les données de tous les tenants coexistent dans
les mêmes tables. Une colonne \texttt{tenant\_id}, indexée, identifie à qui appartient
chaque enregistrement. Toutes les requêtes SQL doivent inclure un filtre
\texttt{WHERE tenant\_id = :id} pour garantir l'isolation logique.

Les avantages sont nombreux~: déploiement simplifié (une seule base de données à
maintenir), migrations de schéma appliquées une seule fois, consommation de ressources
optimisée. Les risques résident dans les bugs de filtrage (oubli du filtre
\texttt{tenant\_id} dans une requête) et dans la moindre isolation en cas d'attaque ou
d'erreur. BCaaS adopte cette stratégie pour son MVP, en s'appuyant sur des mécanismes de
sécurité au niveau applicatif (middleware d'injection automatique du \texttt{tenant\_id})
pour réduire les risques d'oubli.

\subsubsection*{b) Schéma PostgreSQL séparé par tenant}

Chaque tenant dispose de son propre schéma PostgreSQL (un \textit{namespace} de tables)
au sein d'une même base de données. L'isolation est physique au niveau des données mais
les tenants partagent toujours le moteur de base de données et ses ressources.

Cette approche offre une meilleure isolation (il est impossible d'accéder aux données d'un
autre schéma sans changer explicitement de \texttt{search\_path}), des sauvegardes par
tenant plus aisées, et une capacité à personnaliser le schéma par tenant. En contrepartie,
les migrations doivent être appliquées à chaque schéma séparément, ce qui complexifie les
déploiements.

\subsubsection*{c) Base de données séparée par tenant}

Chaque tenant dispose de sa propre instance de base de données. C'est le niveau
d'isolation le plus fort~: une défaillance ou une compromission de la base d'un tenant
n'affecte pas les autres. Cette stratégie est requise dans les secteurs très régulés
(banque, santé) où les données doivent être physiquement séparées.

Le coût opérationnel est significatif~: $N$ tenants impliquent $N$ bases à provisionner,
surveiller, sauvegarder et mettre à jour. Elle est donc rarement justifiée pour un grand
nombre de petits tenants.

% ---- 2.2.3 -------------------------------------------------
\subsection{Analogie avec les réseaux virtuels (VLAN)}
\label{subsec:vlan-analogy}

L'analogie entre le multi-tenant et les VLAN (\textit{Virtual Local Area Networks}) est
particulièrement éclairante. Un VLAN permet de segmenter logiquement un réseau physique
unique en plusieurs réseaux virtuels isolés. Des équipements physiquement connectés au
même commutateur peuvent appartenir à des VLAN différents et ne pas pouvoir communiquer
entre eux, comme s'ils étaient sur des réseaux entièrement séparés.

De même, dans une architecture multi-tenant à schéma partagé, les données de différents
tenants coexistent sur la même infrastructure physique (même serveur, même base de
données) mais sont isolées logiquement par leur \texttt{tenant\_id}, exactement comme les
trames Ethernet sont isolées par leur \textit{tag} VLAN (norme IEEE~802.1Q). BCaaS
exploite cette analogie pour conceptualiser et documenter son architecture.

% ------------------------------------------------------------
\section{Le paradigme \textit{as-a-Service} et le \textit{Business as a Service}}
\label{sec:baas}

% ---- 2.3.1 -------------------------------------------------
\subsection{L'évolution des modèles \textit{as-a-Service}}
\label{subsec:aas-evolution}

Le cloud computing a popularisé une famille de modèles de fourniture de services
informatiques regroupés sous le terme générique \textit{as-a-Service}. L'idée centrale
est de fournir une capacité informatique sous forme de service accessible via une interface
standardisée, déchargeant le consommateur de la responsabilité de construire et d'opérer
l'infrastructure sous-jacente.
\begin{table}[ht]
  \centering
  \caption{Comparaison des modèles \textit{as-a-Service} et positionnement de BCaaS}
  \label{tab:aas}
  \renewcommand{\arraystretch}{1.3}
  \begin{tabularx}{\textwidth}{>{\bfseries}p{2.2cm}
                               >{\raggedright\arraybackslash}X
                               >{\raggedright\arraybackslash}X
                               >{\raggedright\arraybackslash}X}
    \toprule
    \textbf{Modèle} & \textbf{Ce qui est fourni} & \textbf{Exemples concrets} & \textbf{Analogie BCaaS} \\
    \midrule
    IaaS & Infrastructure (serveurs, stockage, réseau) & AWS EC2, Azure VMs, GCP Compute & --- \\
    PaaS & Plateforme d'exécution + middleware & Heroku, Vercel, Render & --- \\
    SaaS & Application complète prête à l'emploi & Salesforce, Gmail, Notion & --- \\
    BaaS (BCaaS) & Briques métier réutilisables via API & BCaaS acteurs, règles, facturation & Couche métier partageable \\
    \bottomrule
  \end{tabularx}
\end{table}
Ces trois modèles forment une progression~: IaaS délègue la gestion du matériel, PaaS y
ajoute celle du système d'exploitation et des middlewares, SaaS va jusqu'à déléguer la
gestion de l'application elle-même. Le consommateur n'a plus qu'à utiliser le service,
sans se préoccuper de ce qui se passe en dessous.

% ---- 2.3.2 -------------------------------------------------
\subsection{Le \textit{Business as a Service} (BaaS)}
\label{subsec:baas-def}

Le \textit{Business as a Service} (BaaS) transpose cette logique au niveau des fonctions
métier. Plutôt que de déléguer de l'infrastructure ou une application, il s'agit d'exposer
comme services réutilisables les briques fonctionnelles communes à tout système
d'information professionnel.

Ces briques récurrentes comprennent typiquement~:
\begin{itemize}
  \item La gestion des acteurs~: utilisateurs, clients, fournisseurs, partenaires et leurs
        rôles respectifs.
  \item La gestion des ressources~: produits, services, inventaires, tarifs.
  \item La facturation et la comptabilité~: devis, commandes, factures, paiements,
        rapprochements.
  \item Les règles métier~: politiques de prix, conditions d'accès, workflows
        d'approbation.
  \item La traçabilité et l'audit~: journaux d'événements, historiques, conformité
        réglementaire.
\end{itemize}

Au lieu que chaque projet réimplémente ces fonctions depuis zéro, le BaaS propose de les
centraliser dans un socle commun consommable via des API. BCaaS est une mise en œuvre
concrète de ce paradigme, conçue comme socle commun à plusieurs instances métier distinctes
pouvant appartenir à des secteurs différents.

% ---- 2.3.3 -------------------------------------------------
\subsection{Avantages et défis du BaaS}
\label{subsec:baas-avantages}

\paragraph{Avantages.}
Réduction drastique du temps de développement des projets consommateurs, homogénéité des
comportements métier entre projets, centralisation de la maintenance et des corrections de
bogues, et mutualisation des coûts d'infrastructure.

\paragraph{Défis.}
Le socle doit être suffisamment générique pour couvrir des cas d'usage variés sans être si
abstrait qu'il devient inutilisable. La gestion des personnalisations par tenant (règles
métier spécifiques à un client) et la compatibilité ascendante de l'API sont des points de
vigilance critiques.

% ------------------------------------------------------------
\section{Architecture hexagonale et \textit{Domain-Driven Design}}
\label{sec:hexa-ddd}

% ---- 2.4.1 -------------------------------------------------
\subsection{L'architecture hexagonale (\textit{Ports \& Adaptateurs})}
\label{subsec:hexagonal}

Proposée par Alistair Cockburn en 2005~\cite{hexagonale}, l'architecture hexagonale 
également appelée \textit{Ports and Adapters} repose sur un principe fondamental~: le
c\oe{}ur métier d'une application ne doit dépendre d'aucun détail technique. La base de
données, le framework web, la messagerie asynchrone, les systèmes externes sont tous des
détails susceptibles de changer, et le domaine métier ne doit pas en être affecté.

L'architecture se structure autour de trois zones~:

\begin{description}[leftmargin=0.5cm]
  \item[\textbf{Le domaine (hexagone central)}]
    Contient la logique métier pure~: les entités, les règles, les calculs. Il est exprimé
    en langage métier et ne dépend d'aucune technologie.

  \item[\textbf{Les ports}]
    Interfaces (au sens de contrats ou d'abstractions) définies par le domaine pour
    interagir avec l'extérieur. Un port entrant (\textit{driving port}) représente ce que
    le domaine offre aux utilisateurs (commandes, requêtes). Un port sortant
    (\textit{driven port}) représente ce dont le domaine a besoin de l'infrastructure
    (persistance, notifications).

  \item[\textbf{Les adaptateurs}]
    Implémentations concrètes des ports. Un adaptateur REST traduit les requêtes HTTP en
    appels de commandes métier. Un adaptateur PostgreSQL traduit les appels de persistance
    en requêtes SQL. Les adaptateurs sont interchangeables sans toucher au domaine.
\end{description}

Les bénéfices de cette architecture sont considérables pour un projet comme BCaaS. Le
c\oe{}ur métier peut être testé indépendamment de toute infrastructure (tests unitaires
rapides, sans base de données ni serveur HTTP). On peut changer de base de données, de
framework web ou de système de messagerie sans réécrire la logique métier. Les adaptateurs
peuvent être développés et remplacés indépendamment.

% ---- 2.4.2 -------------------------------------------------
\subsection{Le \textit{Domain-Driven Design} (DDD)}
\label{subsec:ddd}

Le \textit{Domain-Driven Design}, formalisé par Eric Evans dans son ouvrage fondateur
\textit{Domain-Driven Design: Tackling Complexity in the Heart of Software}
(2003)~\cite{evans2003}, fournit un ensemble de \textit{patterns} et de pratiques pour
modéliser des domaines métier complexes. Il s'articule autour de plusieurs concepts
clés~:

\begin{description}[leftmargin=0.5cm]
  \item[\textbf{Le langage omniprésent} (\textit{Ubiquitous Language})]
    Développeurs et experts métier utilisent exactement le même vocabulaire pour désigner
    les concepts du domaine. Ce langage partagé se retrouve dans le code, les tests, la
    documentation et les discussions. Il élimine les ambiguïtés et les traductions
    coûteuses entre monde technique et monde métier.

  \item[\textbf{Les entités} (\textit{Entities})]
    Objets métier possédant une identité unique et persistante à travers le temps,
    indépendamment de leurs attributs. Un client reste le même client même si son adresse
    change~: c'est l'identité qui compte, pas les valeurs.

  \item[\textbf{Les objets-valeur} (\textit{Value Objects})]
    Objets définis uniquement par leurs attributs, sans identité propre. Une adresse
    postale, un montant en devise, une plage de dates sont des objets-valeur~: deux
    instances identiques sont interchangeables.

  \item[\textbf{Les agrégats} (\textit{Aggregates})]
    Grappes d'entités et d'objets-valeur traitées comme une unité cohérente pour les
    modifications. Chaque agrégat a une racine (\textit{Aggregate Root}) qui est le seul
    point d'accès depuis l'extérieur, garantissant la cohérence des invariants métier.

  \item[\textbf{Les services de domaine} (\textit{Domain Services})]
    Opérations métier qui n'appartiennent naturellement à aucune entité ou objet-valeur,
    telles qu'un calcul de TVA impliquant des règles fiscales complexes.

  \item[\textbf{Les contextes délimités} (\textit{Bounded Contexts})]
    Frontières explicites dans lesquelles un modèle de domaine particulier s'applique. Un
    même mot peut avoir des significations différentes dans des contextes différents~: un
    «~client~» dans le contexte commercial n'est pas le même concept que dans le contexte
    du support technique.
\end{description}

% ---- 2.4.3 -------------------------------------------------
\subsection{DDD et architecture hexagonale : une synergie naturelle}
\label{subsec:ddd-hexa-synergie}

L'architecture hexagonale et le DDD se complètent naturellement. Le DDD fournit le modèle
conceptuel du domaine (ce qu'on modélise et pourquoi), tandis que l'architecture
hexagonale fournit la structure technique pour implémenter ce modèle de manière isolée et
testable (comment on l'organise dans le code).

Ensemble, ils permettent de construire un \textbf{monolithe modulaire}~: une application
unique et cohérente, mais découpée en modules aux frontières explicites (les
\textit{Bounded Contexts}) et aux interfaces bien définies (les ports). Cette structure
prépare une éventuelle migration vers les microservices~: chaque module peut, à terme,
être extrait et déployé indépendamment sans réécriture du domaine, car les dépendances
inter-modules passent déjà par des interfaces, non par des appels directs.

% ------------------------------------------------------------
\section{Positionnement de BCaaS}
\label{sec:positionnement}

% ---- 2.5.1 -------------------------------------------------
\subsection{Ce que la littérature fournit déjà}
\label{subsec:litterature-existante}

La revue de la littérature met en évidence que chacun des éléments conceptuels mobilisés
par BCaaS existe et est documenté indépendamment~:

\begin{itemize}
  \item Les modèles en couches OSI et TCP/IP sont des standards établis depuis plusieurs
        décennies, avec une littérature abondante et des implémentations universelles.
  \item Les stratégies multi-tenant sont documentées dans la littérature académique et dans
        les guides de bonnes pratiques des fournisseurs cloud (Microsoft Azure, AWS).
  \item Le paradigme \textit{as-a-Service} est largement décrit et analysé, notamment dans
        les travaux du NIST sur le cloud computing~\cite{nist}.
  \item L'architecture hexagonale et le DDD sont des méthodologies reconnues, avec des
        communautés actives et des implémentations de référence dans de nombreux langages.
\end{itemize}

% ---- 2.5.2 -------------------------------------------------
\subsection{Le \textit{gap} identifié}
\label{subsec:gap}

Cependant, la littérature existante traite ces éléments séparément, dans des silos
disciplinaires distincts~: les modèles réseau dans la littérature des systèmes distribués,
le multi-tenant dans la littérature SaaS, l'architecture hexagonale dans la littérature du
génie logiciel. Il n'existe pas, à notre connaissance, de travail proposant une grille de
lecture unifiée qui utilise les métaphores réseau pour conceptualiser et documenter un
système logiciel multi-tenant à architecture hexagonale.

% ---- 2.5.3 -------------------------------------------------
\subsection{La contribution de BCaaS}
\label{subsec:contribution}

L'originalité de BCaaS ne réside donc pas dans l'invention de l'un de ces éléments, mais
dans leur \textbf{unification sous une grille de lecture réseau cohérente et systématique}.
Cette grille se manifeste à travers plusieurs analogies structurantes~:

\begin{description}[leftmargin=0.5cm]
  \item[\textbf{La pile métier comme pile protocolaire}]
    Les couches de l'application (présentation, application, domaine, infrastructure) sont
    interprétées comme les couches d'une pile réseau, avec des responsabilités distinctes
    et des interfaces explicites entre couches.

  \item[\textbf{L'isolation multi-tenant comme VLAN}]
    Le partage d'infrastructure avec étiquetage logique des flux par \texttt{tenant\_id}
    est l'exact équivalent logiciel du marquage de trames par \textit{tag} VLAN dans les
    réseaux commutés (IEEE~802.1Q).

  \item[\textbf{La facture comme accusé de réception}]
    Le cycle de vie d'une facture (émission, accusé, paiement, rapprochement) est modélisé
    sur le \textit{pattern} ACK/NACK des protocoles réseau fiables~: chaque étape confirme
    la bonne réception et le traitement du message précédent.

  \item[\textbf{L'administration comme plan de contrôle SDN}]
    La séparation entre le plan de données (traitement des transactions métier) et le plan
    de contrôle (configuration des règles, gestion des tenants, supervision) reflète
    l'architecture \textit{Software-Defined Networking}, où le plan de contrôle est
    découplé du plan de données pour plus de flexibilité.
\end{description}

Cette cohérence conceptuelle (utiliser un seul cadre métaphorique pour relier des
concepts habituellement traités séparément) constitue l'apport central du présent
travail. Elle offre un langage commun accessible aussi bien aux profils techniques
(développeurs, architectes) qu'aux profils familiers des réseaux, et facilite la
communication entre équipes aux expertises différentes.

La transposition de ces métaphores en décisions d'architecture concrètes sera détaillée au
chapitre~\ref{chap:archi}, qui présente le modèle conceptuel de BCaaS.

% ------------------------------------------------------------
\section*{Synthèse}
\addcontentsline{toc}{section}{Conclusion du chapitre}
\label{sec:conclusion-chap2}

Ce chapitre a passé en revue les quatre piliers conceptuels sur lesquels repose BCaaS~:
les modèles en couches OSI/TCP-IP, les architectures multi-tenant, le paradigme
\textit{as-a-Service} étendu au niveau métier, et l'architecture hexagonale couplée au
\textit{Domain-Driven Design}. Pour chacun de ces piliers, nous avons présenté les
concepts fondamentaux, leurs avantages et leurs limites, ainsi que leurs applications
concrètes dans le contexte de BCaaS.

La section~\ref{sec:positionnement} a montré que l'originalité de BCaaS tient à la
combinaison et à l'unification de ces approches sous une grille de lecture réseau
cohérente une perspective que la littérature n'exploite pas jusqu'ici de manière
systématique. Le chapitre suivant s'appuiera sur cet état de l'art pour présenter le
modèle d'architecture détaillé de BCaaS et expliciter comment chaque analogie réseau se
traduit en décisions de conception concrètes.